% Editorial correspondence: gestion/mapa_transportes_auxiliares.md.
\section{Operational memory and invariants of an evolving structure}
\label{md:lectin:seccion}

\subsection{Observable value, capacity to act, and compatibility}
\label{md:lectin:valor}

To know an object is to determine what its observable value preserves and which
aspects of its capacity to act require an additional description.
In the HMT construction, two present readouts may coincide while
the orientation and memory retained produce different responses
to the same continuation. Genealogy acquires
physical and mathematical content when it determines these differences in response.

The origin of this investigation is the APP--TRIT--TPK composition, its
enriched state, and the joint discrete structure. The coordinates
$\pi,\varphi,e,\alpha$ retain their role as correlated outputs.
The capacity, reciprocity, and action readout operators intervene afterwards,
with their own domains. The following compositions examine those
readout operators and retain the generator as their antecedent.

One conceptual characterization of constants is to study them
as compatibility invariants of an evolving structure.
For each constant, this characterization requires specifying the
transformations under which it remains invariant, the information
that may vary, and the readout operator through which it is identified. Radial
reciprocity provides a proved case: the action section
remains invariant when the radius is inverted and the
corresponding change of orientation is retained. Extending this
invariance to the full holonomy requires the explicit composition
of the operations involved.

\subsection{Retained memory and operationally visible memory}
\label{md:lectin:memoria-visible}

In the joint representation of readout and memory, the readout is updated
by $x'=Tx+Dm$. For two states with the same
present readout, one obtains
\begin{equation}
 \Delta x'=D\,\Delta m.
 \label{md:lectin:diferencia-memoria}
\end{equation}
The genealogical difference is visible at that step precisely when
$D\Delta m\ne0$. The identity identifies the continuation that
transforms a difference in memory into a difference in readout.

Conversely, for a stationary transport $C$ in
this representation, one has
\begin{equation}
 D=\frac{\sqrt8}{9}(C-I),\qquad R=\frac{I+8C}{9}.
 \label{md:lectin:bloques-memoria}
\end{equation}
If $C\Delta m=\Delta m$, then $D\Delta m=0$ and
$R\Delta m=\Delta m$. This difference in memory remains invisible
under repetitions of the same protocol. A different operation may
transfer it to the readout.

The sufficiency of a description therefore concerns the
differences it must preserve in order to predict responses to the
admissible operations. The full genealogy may contain more
information than a particular observation requires. The
observer's scope determines which reduction retains predictive
capacity. This reduction is studied from the transports produced
by HMT and preserves the provenance of its states and operators.

In the stationary case, the invisible memory differences are
exactly $\operatorname{Fix}(C)$. For a collection of available operations,
the differences invisible under every continuation
form the intersection of their fixed spaces. Necessity follows
by taking each operation as the first step. Sufficiency is proved
by induction: every operation fixes the difference in question and
its transfer block annihilates that difference, so that it never
appears in the readout. The result concerns this linear
representation; the chronological counter of the holonomy retains its own map
into it.

\subsection{Regularity of operations and prolongation of the state}
\label{md:lectin:regularidad}

Vacancy cadences distinguish several levels of closure. A
TRIT regime may return while the capacity phase remains
displaced. A nonadic phase may return while the
turn counter increases. Repetition of a type of operation is
compatible with the production of a new history.

This distinction makes it possible to study the regularity of an
unbounded prolongation: a finite rule may order the evolution without every
new state reproducing an earlier one. Finiteness of the rule and
finiteness of the set of reachable states are different properties.
An aperiodic decimal readout may arise from dynamics subject
to strict relations.

In the correlated regions, the object of study is the relation
that jointly determines the next block of each region, the
properties preserved by the mirror, and the memory that makes the
prolongations compatible. The $21/22$ and $6/7$ cadences are capacity
observations; their concrete maps relate them to the regional
readout operators. The proof of arbitrary depth and the explanation of
the shared structure are distinct and
complementary results.

Composition of the cadences separates the levels exactly:
completing the three changes of the selector leaves a capacity
displacement of three or six classes modulo nine. The durations $437$
and $459$ quantify this distinction in the corresponding construction;
their role here is to describe those cadences, without promoting them to
additional fundamental constants.

\subsection{Action character and oriented radial reconstruction}
\label{md:lectin:h5}

The character $H_5$ and the regional continuation determine an action
section. The radial form makes it possible to reconstruct that section when
orientation is retained. At fixed $\alpha$, the coordinate
$\chi_{\mathrm{rad}}=\varepsilon\log q_\varepsilon$ classifies
reciprocity and is monotonically related to the return section:
\begin{equation}
 \frac{\eta_{\mathrm{ret}}}{\alpha}
 =-500\tanh\!\left(\frac{\chi_{\mathrm{rad}}}{2}\right)-1.
 \label{md:lectin:accion-radial}
\end{equation}
The oriented radial form and the action section are subject to
a joint constraint. Action incorporates a condition of
geometric compatibility; its coefficients, regional correction,
and orientation must remain coordinated.

This relation distinguishes two levels of knowledge.
Reparametrization proves equivalence of descriptions and makes it possible
to recover either from the other. A comparison between readout operators provides
additional information when both are evaluated from the preceding structure
without introducing the value being checked as an input.
Computing the radius through the inverse of the action formula
and substituting it back verifies an algebraic identity.
Producing the geometric readout through its antecedent operations
and comparing it with the action character tests the causal
composition. Each check has its own role.

The distinction identifies the developments that must be assembled and
executed; it retains HMT genealogy as their common root. The comparison
must be independent of the value being checked, not of
the shared origin of its readout operators.

\subsection{Action, information, and temperature}
\label{md:lectin:termica}

The state of the admissible continuations and their weights determines
dimensionless information. The transducer $k_B$ realizes its readout in
entropy units. Action and period enter the relation
\begin{equation}
 k_B\Theta_{\mathrm{clk}}\log3=\frac{h}{T_{108}}.
 \label{md:lectin:accion-informacion}
\end{equation}
The three equiprobable alternatives produce $\log3$; a nonuniform
distribution retains its weights in the evaluation of information.

The composition makes it possible to follow how the same evolution changes
operational memory, compatible continuations, and the energy
readout. The length of a history and its temperature correspond
to different readouts. Thermal evaluation requires identifying the
statistical state and the exchange operation that realize it.
The precise problem is to determine which generated transformations
preserve action, which redistribute accessible information,
and which produce a calculable thermal difference.

When a path closes in a projection while its memory changes,
evaluating its work requires declaring the space in which
the cycle is considered complete. This is the problem of
conservativity. Noncommutativity, by contrast, characterizes
sensitivity to the order of operations. Preserving the distinction
makes it possible to study the relation between the two properties.

\subsection{Comparison programme and hierarchy of checks}
\label{md:lectin:contrastes}

\paragraph{Histories with a common readout and different future responses.}
The first comparison brings together two admissible routes produced by TPK
that coincide under one readout operator, retains their memories, and applies the same
continuation. The quantity examined is the exact separation of their
subsequent readouts and the memory component that determines it.
The difference between value and structure is thus expressed through a
difference in predictive capacity. The examples already constructed
retain their provenance and can be prolonged within this scheme.

\paragraph{Compatibility between regional action and geometric form.}
The second comparison follows the $H_5$ readout operator and its continuation, and
the geometric readout operator with its own antecedents. Their compatibility
is examined before and after a return, distinguishing the limiting
section from its finite approximations. The comparison identifies what
remains constant while memory advances. The proved radial invariance
provides the reference case and retains its specific
scope.

\paragraph{Order of operations and memory re-entry.}
The third comparison uses two effective TPK operations and
retains the intermediate charts. The $1:8$ partition of the auxiliary
realization corresponds to the regime of a fibre with identity reference.
The application must verify that regime or transport the corresponding
charts. The mathematical experiment compares the two orders
with and without memory re-entry; the physical realization must specify
the preparation, interaction, and observable that represent those
operations.

The first two comparisons establish the conceptual priority: they determine
which information has dynamical efficacy and which invariants support
the compatibility of the architecture. The third provides a
direct quantitative response within its proved regime.
This hierarchy preserves the auxiliary role of the energy realization.

\subsection{Scientific comparison and scope of interpretation}
\label{md:lectin:comparacion}

Memory in reduced dynamics has precedents in mathematical
physics. The Mori--Zwanzig formalism makes it possible to study the
evolution and memory terms of reduced observables; an
example is the analysis by Hudson and Li.\footnote{Hudson and Li,
\emph{Coarse-graining of overdamped Langevin dynamics via the
Mori--Zwanzig formalism},
\href{https://arxiv.org/abs/1810.08175}{arXiv:1810.08175}.}
The relevant comparison concerns the concrete composition examined
here: APP--TRIT--TPK origin, correlated regions, action,
orientation, and reconstruction. Mentioning memory alone does not
determine priority for that composition.

The expository contribution consists in specifying consequences and
comparisons within the construction. A historical assessment of
novelty and physical corroboration retain their specific
tests. The mathematical programme establishes objects and operations
that can be subjected to scientific comparison without anticipating the
outcome of those assessments.

Stored memory and operationally visible memory are characterized
here through the kernel of $D$ and the fixed space of
$C$. This characterization follows from the declared update
and retains the scope of its representation. The three
comparisons in Section~\ref{md:lectin:contrastes} constitute a
work programme; their formulation distinguishes the definition of the
mathematical experiment, the execution of TPK routes, and the realization
of a physical experiment.
